\scanpage{81}
% archive foliation: 83
\[
\left\{\begin{gathered}
\text{paires de pantalons }\deleted{\text{\unsure{dég.}}}\inserted{\text{\unsure{généralisés}}}\inserted{X',X''},\ \inserted{\text{\unsure{à trous}}}\ \text{indexés}\\
\text{par }\varepsilon'\amalg\{s\},\quad\varepsilon''\amalg\{t\},\quad\text{avec }s,t\ \text{correspondant}\\
\text{aux }\deleted{\text{bords}}\ \text{composantes ponctuelles du}\\
\text{bord, avec les paramètres }(\lambda_i)_{i\in\varepsilon'\amalg\varepsilon''}\in\unsure{(\mathbb R^{*+})^{\varepsilon'\amalg\varepsilon''}}\\
\text{et un isom d'identification}\\
\text{des }\mathbb C^*\wedge_{\pm1}(\varepsilon'\wedge\varepsilon'')
\end{gathered}\right.
\]
Une fois ce cas compris, et revenant
alors au cas général, il faudrait déjà,
pour un type de découpage donné
incarné par un diagramme MD
dans $md_{g,\nu}$, comprendre l'application
sur les groupes fondamentaux
induite par
\diagram{\begin{tikzcd}[ampersand replacement=\&]
\unsure{\varphi})\quad MD\deleted{B}_{\unsure{g}} \arrow[r] \& M_{\mathrm{st}}N\deleted{\unsure{D}}_{g,\nu},
\end{tikzcd}}
ceux du deuxième membre étant
essentiellement compris, en termes de
groupes de Teichmüller spéciaux
--- \uline{mais il doit en être pareil pour le
premier} ! Également, il ne devrait
pas être bien dur de vérifier si $\unsure{\varphi}$
